\section{What is asserted in intuitionistic primitive recursive arithmetic}
\label{sec:foundations}

The natural numbers are the only object-theoretic domain. We use the
ordinary first-order presentation of $\ipra$: a symbol for each primitive
recursive definition; its defining equations; $Sx\ne0$ and
$Sx=Sy\Rightarrow x=y$; intuitionistic predicate logic with equality; and
induction for quantifier-free formulas. Concretely, a primitive recursive
symbol has an arity and is generated from zero, successor, projections,
composition, and
\[
 R(f,g)(0,\mathbf x)=f(\mathbf x),\qquad
 R(f,g)(n+1,\mathbf x)=g(n,R(f,g)(n,\mathbf x),\mathbf x).
\]
Induction is
\[
 A(0,\mathbf u)\wedge
 \forall n\,[A(n,\mathbf u)\Rightarrow A(n+1,\mathbf u)]
 \ \Rightarrow\ \forall n\,A(n,\mathbf u),
 \qquad A\text{ quantifier-free}.
\]
The logical rules are hypothesis, substitution of terms, equality
reflexivity and Leibniz substitution, explosion from $\bot$, introduction
and elimination for $\wedge,\vee,\to,\forall,\exists$, with the usual
fresh-variable conditions on $\forall$-introduction and
$\exists$-elimination. Neither excluded middle for arbitrary formulas nor
double-negation elimination is added. Decidable comparisons of integer
codes do admit case analysis, proved by their defining equations.

This is precisely the intended arithmetic interface of
\texttt{ConciseIPRA/PR.lean} and \texttt{Deduction.lean} at the repository
revision recorded in Section~\ref{sec:lean}. The inspected theorem
\texttt{pra\_iff\_native} transfers a completed derivation to the native
concise calculus. It does not automatically produce a derivation of an
analytic statement~\cite{ProveIt}.

\subsection{Arithmetic meaning of real and complex notation}
An integer is a sign and a magnitude; a rational is an integer numerator
and a positive natural denominator. Reduction to lowest terms is optional.
Rational equality and order use cross multiplication. Pairing, finite
lists, bounded search, factorials, binomial coefficients, and finite sums
and products are primitive recursive. For example a bounded universal
statement about a decidable predicate $P$ is the equation
$\prod_{j<b}\chi_P(j)=1$. Thus bounded induction invariants can be
written without bounded quantifiers.

The elementary factorization facts used later have the same finite
status. Trial division is bounded by the integer being factored, and
its loop invariant is ``the original integer is the product of the
recorded factors and the current remainder.'' Extended Euclid can be
run for at most $a+b+1$ steps on positive inputs $a,b$, leaving a fixed
state once a remainder is zero. Its invariant records both the gcd and
the integer linear combination for each remainder. Remainders strictly
decrease until zero. If a prime $p$ does not divide $a$, its gcd with
$a$ is one, so the recorded Bezout equality proves that $p\mid ab$
implies $p\mid b$. Repeated application gives uniqueness of the finite
factorization. The sequence entries, coefficient bounds obtained by
iterating the elementary update, and loop invariants are all coded
primitive recursively. This supplies the finite arithmetic behind the
Euler product and divisor identities.

A \emph{name} is a fixed primitive recursive rational evaluator
$a(\mathbf u,k)$, together with arithmetic hypotheses or proofs of
\[
 |a(\mathbf u,k)-a(\mathbf u,l)|\le2^{-k}+2^{-l}.
 \tag{F1}\label{eq:regular}
\]
Complex names use pairs of rationals and the norm
$\|x+iy\|_1=|x|+|y|$. We omit $\mathbf u$ from the notation.
Two names are equal when
\[
 \forall k,l\quad \|a(k)-b(l)\|_1\le2^{-k}+2^{-l}.
 \tag{F2}\label{eq:name-equality}
\]
The name $a$ is real and at least $d\in\Q$ when
$a(k)\ge d-2^{-k}$ for every $k$. Complex apartness is a supplied rational
$d>0$ with $|a|\ge d$, expressed by rational squared-modulus inequalities
with the corresponding approximation error. A convenient equivalent
implementation works with rectangular rational balls.

The Euclidean modulus in estimates has
$|z|\le\|z\|_1\le2|z|$. To implement a Euclidean error bound as a
$1$-norm bound, multiply it by two, or decrease the error budget by one
bit. Squared-modulus comparisons avoid square roots altogether.

These definitions involve only number quantifiers after the fixed
functions have been named. They do \emph{not} quantify over all functions
$\N\to\Q$, or assert a total primitive recursive evaluator for arbitrary
primitive recursive program codes. Such a universal evaluator would give
the diagonal primitive recursive function $n\mapsto U(n,n)+1$, an
immediate contradiction. Statements for general named inputs are therefore
\emph{schemas in a fixed evaluator}, whereas statements for rational
inputs and numerical parameters are single uniform first-order statements.

To check transitivity of (F2), compare $a(k)$ to $c(l)$ through $b(j)$.
The upper bound is $2^{-k}+2^{-l}+2^{1-j}$. A positive rational gap from
$2^{-k}+2^{-l}$, if present, has a dyadic below it found from its numerator
and denominator lengths. Choosing that $j$ contradicts the bound.
Only a rational comparison is used. The same argument removes every
``arbitrarily small error'' in this article.

\subsection{The finite-precision proof discipline}
All limiting notation below abbreviates a construction. For example, if
finite expressions $F_N$ obey
\[
 \|F_M-F_N\|_1\le C/N\quad(M\ge N\ge N_0),
\]
choose $N(k)=N_0+(1+\lceil C\rceil)2^{k+4}$ and evaluate $F_{N(k)}$
within $2^{-k-4}$. A fixed shift gives a regular name. There is no appeal
to unrestricted Cauchy completeness or monotone convergence.

The elementary error rules are
\begin{align}
 \left\|\sum_{j<v}a_j-\sum_{j<v}b_j\right\|_1&\le v\epsilon,
       \label{eq:f-sum}\\
 \left\|\prod_{j<v}a_j-\prod_{j<v}b_j\right\|_1
       &\le vM^{v-1}\epsilon
       \quad(\|a_j\|_1,\|b_j\|_1\le M),\label{eq:f-prod}\\
 |a^{-1}-b^{-1}|&\le2\epsilon/d^2
       \quad(|a|\ge d,\ |a-b|\le\epsilon\le d/2).
       \label{eq:f-inv}
\end{align}
The second follows by replacing one factor at a time; the third follows
from $(b-a)/(ab)$. In the $1$-norm, harmless additional factors of two
are included in the precision budget. Integer majorants suffice. An error
coefficient $C\ge1$ can always be absorbed by the deliberately excessive
precision shift $C+k+5$, since $C2^{-C-k-5}<2^{-k-4}$.

\begin{proposition}[Arithmetic reading of the proofs]\label{prop:arithmetic-reading}
For each displayed fixed family of input names, the finite estimates and
algorithms in this article are proofs in $\ipra$ under their stated name
and domain certificates. The quantifiers over precision are introduced
only after the finite rational calculation has been proved.
\end{proposition}
\begin{proof}
Fix the desired precision $k$, all truncation lengths, and all numerical
domain parameters. Expand rational operations into operations on integer
codes. Use the displayed bounds to assign a smaller rational error to
each of the finitely many evaluations. The invariant in a sum, product,
recursion, or mesh calculation is then a finite rational identity or
inequality, equivalently an equation for a primitive recursive Boolean
checker. Apply quantifier-free induction to that invariant. If a finite
list of assertions is needed, use a running conjunction coded by a
Boolean product. Hypotheses on input names are instantiated only at the
finitely many stages used by this calculation. Introduce $k$ universally
at the end and apply (F2) or its inequality version.

This argument is a translation of the particular finite proofs given
below. It is not the false general principle that every true statement
about primitive recursive functions is provable in $\ipra$. In particular,
we never induct on an unbounded-precision assertion such as ``these two
names are equal.'' The quantitative continuation proof later explicitly
keeps its precision losses inside a finite induction.
\end{proof}

\subsection{Elementary functions}
For an integer $B\ge1$, $|z|\le B$, and $m+1\ge2B$, the finite
exponential polynomial satisfies
\[
 \left|e^z-\sum_{j=0}^m\frac{z^j}{j!}\right|
 \le\frac{2B^{m+1}}{(m+1)!}.                     \tag{F3}\label{eq:exp-tail}
\]
Here the left side is initially notation for the difference of two finite
polynomials, with the same bound uniformly in the larger length. The
ratios beyond degree $m$ are at most $1/2$. The sum of the absolute terms
at radius $B$ is at most $3^B$: first $\sum1/j!<3$ follows from the
factorial recurrence, then finite multinomial multiplication gives the
integer-$B$ bound. Thus $m=4B+k+6$ suffices for error $2^{-k}$.
Finite binomial convolution, with (F3), proves $e^{z+w}=e^ze^w$.
For $|z|\le B$, $e^z e^{-z}=1$ gives $|e^z|\ge3^{-B}$.

For positive $q$ define, with $v=(q-1)/(q+1)$,
\[
 \log q=2\sum_{j\ge0}\frac{v^{2j+1}}{2j+1},\qquad
 |\text{tail after }L\text{ terms}|
 \le\frac{2|v|^{2L+1}}{(2L+1)(1-v^2)}.          \tag{F4}\label{eq:log-tail}
\]
On $B^{-1}\le q\le B$, $|v|\le1-1/B$, so
$L=B(k+2B+4)$ suffices. The inequality
$(1-1/B)^B\le1/2$ for $B>1$ follows from
$(1+1/(B-1))^{B-1}\ge2$, with an extra factor if necessary; the case
$B=1$ is immediate. Alternatively binary range reduction gives the
faster algorithm in the code. Formal differentiation of the finite sums
and their geometric derivative tails proves $(\log q)'=1/q$.
The addition law for the logarithm follows by differentiating
$\log(ab)-\log a$ in $a$ and evaluating at $a=1$, with the finite
fundamental theorem below. Thus positive-real powers mean
\[
 x^s=\exp(s\log x),\qquad x>0.
\]
All powers in this article use this convention; a global complex logarithm
is unnecessary.
For explicit verification, the $L$-term logarithm polynomial satisfies
$\ell_L'(q)=(1-v^{2L})/q$, so its derivative error is at most
$B|v|^{2L}$ on $[B^{-1},B]$. Differentiating
$\exp(\log q)/q$ gives zero and its value at one is one; the finite
fundamental theorem therefore proves $\exp(\log q)=q$.
Conjugation and the exponential product law give
$|e^z|^2=e^{2\Re z}$. For real $x$, $e^x=(e^{x/2})^2$ is positive
by its explicit nonzero lower bound, so $|e^z|=e^{\Re z}$.


Sine is the odd factorial series and cosine is the even factorial series. Their addition
laws follow from the exponential product identity. Put
\[
 A(u)=\int_0^u\frac{dt}{1+t^2},\qquad \pi=4A(1).
\]
Integrating a finite geometric identity gives the alternating arctangent
sum with error $1/(2L+1)$. Thus $\pi$ has a PR name, and finite rational
bounds give $3<\pi<4$. Direct differentiation shows
\[
 e^{2iA(u)}=\frac{1+iu}{1-iu},\quad
 \sin(2A(u))=\frac{2u}{1+u^2},\quad
 \cos(2A(u))=\frac{1-u^2}{1+u^2}.                \tag{F5}\label{eq:trig-param}
\]
For example the derivative of $e^{-2iA(u)}(1+iu)/(1-iu)$ is zero
and its initial value is one. At $u=1$ this gives $e^{i\pi/2}=i$.
On $0\le u\le1$, $(2A)'\ge1$. Its inverse is constructed by bounded
approximate-sign bisection: on a bracket of width $w$, evaluate the
midpoint discrepancy within $w/16$; a discrepancy outside $[-w/8,w/8]$
selects a half bracket, and otherwise the inverse is within $w/4$ of the
midpoint. The bracket width halves in either case. This proves sine
positivity on the first quadrant without an unrestricted intermediate
value theorem. Since $\sin''=-\sin$, finite mesh comparison with the
endpoint chord gives
\[
 \sin t\ge 2t/\pi\quad(0\le t\le\pi/2).       \tag{F6}\label{eq:sine-chord}
\]
One proof of the chord comparison sums the fact that successive secant
slopes decrease, itself obtained by integrating the nonpositive second
derivative on a finite mesh. Periodicity and reflection extend the
bounds to the other quadrants.

\subsection{Finite calculus and local analyticity}
For a function given by fixed rational evaluations on a rational interval
$[a,b]$, with supplied Lipschitz bound $L$, use the left mesh sum
\[
 I_M=\frac{b-a}{M}\sum_{j=0}^{M-1}
 f\!\left(a+\frac{j(b-a)}M\right).
\]
Comparison on a common refinement gives
\[
 |I_M-I_{M'}|\le L(b-a)^2(1/M+1/M').             \tag{F7}\label{eq:quad}
\]
An evaluation error $\epsilon$ per summand adds $(b-a)\epsilon$.
Consequently this is an actual PR integral name. On finitely many smooth
cells, apply the estimate on each cell separately; discontinuities at
integer cell endpoints cause no difficulty.

Suppose a supplied derivative $g$, with a supplied Lipschitz bound on the
interval, obeys the local estimate
\[
 |f(x+h)-f(x)-hg(x)|\le |h|\omega(|h|),
\]
with an explicit PR modulus $\omega$. Summing over the mesh telescopes,
so $|f(b)-f(a)-I_M(g)|\le(b-a)\omega((b-a)/M)$. This is the
fundamental theorem needed here. For powers and exponentials the local
estimate follows by subtracting the first two terms of their factorial
series; for rational functions it follows by multiplying out the exact
rational difference. Product differentiation and integration by parts
follow from their finite difference identity, whose additional defect is
bounded by the product of the increments. Repeated integration by parts
therefore means a finite repetition with its error allocated in advance.
Changes of variable used below follow by comparing images of mesh cells,
using the explicit derivative and second-order remainder of the given
substitution. Double integrals are finite double sums with uniform error
bounds before the order of summation is changed.

An assertion of holomorphy means more than pointwise evaluation: on each
specified disk there are PR Taylor coefficient names and a bound
$\sum |a_j|r^j\le M$ at a larger radius $r$. On $|h|\le r/2$ the
Taylor tail is at most $M2^{-d-1}$ after degree $d$, and
\[
 |f(c+h)-f(c)-a_1h|\le M|h|^2/r^2.
\]
These are finite geometric estimates. Products use finite convolution;
translation of center uses the binomial formula on a larger disk.
Inversion uses an apartness witness $d$ and a disk on which the
nonconstant part has size at most $d/2$, followed by a geometric series.
These rules, not an unqualified uniform-limit theorem, justify all
analytic language in the sequel.
